\subsection{Synthetic Division and Long Division of Polynomials}

\begin{definition}[Long Division]
  An algebra method used to one polynomial by another of the same or a lower degree. It works just like a regular number long division.
\end{definition}

\begin{definition}[Synthetic Division]
  A fast, streamlined shortcut for dividing a polynomial by a linear factor of the form $x - k$. It replaces traditional long division by using only the numerical coefficients.
\end{definition}

\textbf{Note: } \textit{Factor Theorem:} ``If upon division you get a remainder of $0$, your division is a factor of your dividend.''

\begin{example}[Long Division]
  \hfill \\
  $f(x) = 2x^3 - x^2 + 2x - 3 ; x = 1$ \\
  $f(1) = 2(1)^3 - (1)^2 + 2(1) - 3 $ \\
  $f(1) = 0$ \\
  $x - 1 = 0$ \\
  \polylongdiv{2x^2+x-3}{x-1} \\
  $\frac{2x^3}{x} = 2x^2 \quad$
  $\frac{x^2}{x} = x \quad$
  $\frac{3x}{x} = 3$ \\
  $f(x) = (x - 1)(2x^2 + x + 3)$
\end{example}

\begin{example}[Synthetic Division]
  \hfill \\
  $
  \begin{array}{c|rrrr}
    1 & 2 & -1 & 2 & -3 \\
    & \downarrow & ^+2 & ^+1 & ^+3 \\
    \hline
    & 2 & 1 & 3 & 0 
  \end{array}
  $ \\
  $2x^2 + x + 3$ \\
  $f(x) = (x - 1)(2x^2 + x + 3)$ \\
\end{example}

\begin{example}[Other Examples]
  $$g(x) - 3x^6 + 82x^3 + 27 ; x = -3$$
  $g(-3) = 0$
  \begin{enumerate}
  \item Long Division: \\
    \polylongdiv{3x^6+82x^3+27}{x+3} \\
    $g(x) = (x + 3)(3x^5 - 9x^4 + 27x^3 + x^2 - 3x + 9)$
  \item Synthetic Division: \\
    $
    \begin{array}{c|rrrrrrr}
      -3 & 3 & 0 & 0 & 82 & 0 & 0 & 27 \\
      & \downarrow & ^+-9 & ^+27 & ^+-81 & ^+-3 & ^+9 & ^+-27 \\
      \hline
      & 3 & -9 & 27 & 1 & -3 & 9 & 0 \\
    \end{array}
    $ \\
    $g(x) = (x + 3) (3x^5 - 9x^4 + 27x^3 + x^2 - 3x + 9)$
  \end{enumerate}
\end{example}

\begin{example}[Example with a Non Zero Remainder]
  \hfill \\
  $$h(x) = 4x^3 - 3x^2 - 8x + 4 ; x = 2$$ \\
  $h(2) = 8$
  \begin{enumerate}
  \item Long Division: \\
    \polylongdiv{4x^3-3x^2-8x+4}{x-2} \\
    $4x^2 + 5x + 2 + \frac{8}{x - 2}$ \\
    $h(x) = (x - 2)(4x^2 + 5x + 2 + \frac{8}{x - 2})$
  \item Synthetic Division: \\
    $
    \begin{array}{c|rrrr}
      2 & 4 & -3 & -8 & 4 \\
      & \downarrow ^+8 & ^+10 & ^+4 \\
      \hline
      & 4 & 5 & 2 & 8 \\
    \end{array}
    $ \\
    $4x^2 + 5x + 2 + \frac{8}{x - 2}$ \\
    $h(x) = (x - 2)(4x^2 + 5x + 2 + \frac{8}{x - 2})$
  \end{enumerate}
\end{example}
